\section{Introduction}
\label{sec:intro}

Large language model (LLM) agents now carry out multi-step research tasks rather than
answering single questions. An agent interleaves reasoning with tool calls---searching,
reading files, running code---and observes the results before choosing its next
action~\cite{yao2023react,wang2024survey,xi2025rise}. Researchers have adopted these tools
quickly. In a 2025 survey of 2,430 researchers, 84\% reported using AI tools, up from
57\% a year earlier~\cite{wiley2025explanaitions}; earlier surveys and corpus studies
found the same trend in attitudes and in the text of published
papers~\cite{vannoorden2023ai,liang2024mapping}. Electrical and computer engineering is
no exception: LLMs already generate and debug register-transfer-level hardware
descriptions~\cite{thakur2023benchmarking,blocklove2023chipchat,liu2023verilogeval}, and
domain-adapted models assist chip-design engineers~\cite{liu2023chipnemo}.

Practice has not kept pace with adoption. In the same survey, 57\% of researchers named a
lack of guidelines and training as a barrier, and the share concerned about inaccurate
output rose from 51\% to 64\%~\cite{wiley2025explanaitions}. The concern is well
founded. Language models produce fluent but false statements~\cite{ji2023survey};
in one study, 55\% of the bibliographic citations produced by GPT-3.5 and 18\% of those
produced by GPT-4 were fabricated~\cite{walters2023fabrication}. Models tend to agree with
their users~\cite{sharma2023sycophancy}, coding agents can satisfy tests by exploiting
them instead of meeting the specification~\cite{zhong2025impossiblebench}, and an
analysis of more than 1,600 execution traces from seven multi-agent frameworks
identified 14 distinct failure modes~\cite{cemri2025multiagent}. Meanwhile, venues and
funders agree that an AI tool cannot be an author and that the human authors take full
responsibility for everything a paper
contains~\cite{arxiv2025genai,thorp2023chatgpt,nature2023tools,ec2025living}. The
researcher therefore remains accountable for work that is increasingly performed by
systems that act on their own.

Two bodies of work bracket this problem without solving it. Multi-agent frameworks and
``AI scientist'' systems coordinate agents with each other and push toward greater
autonomy~\cite{wu2023autogen,hong2023metagpt,qian2024chatdev,lu2024aiscientist,schmidgall2025agentlab};
the human, where present, is an approver of last resort, and work passes between agents
as conversation. Research on levels of automation and human supervisory control models
exactly the quantity that matters---how much a human operator can oversee---but it was
developed for operators of vehicles and robots, not for agents that deliver asynchronous
documents and code whose cost lies in
verification~\cite{sheridan1978human,parasuraman2000model,crandall2005validating,cummings2008predicting}.
Recent proposals grade the autonomy of AI agents by the role the user
plays~\cite{feng2025levels}, but they do not say how to assign a level to a given task or
how many agents a single person can responsibly supervise. What is missing is an
operational framework that tells a researcher three things: how much autonomy to grant
each task, how work should pass between agents so that it stays reviewable, and how
large a team their attention can support.

This paper develops such a framework and examines it against three cases. It condenses
and formalizes the central argument of the author's open textbook \textit{AI for EE
Research}~\cite{mccreary2026aiforee}: that the researcher's role shifts from performing
each task to \textit{orchestrating} a team of planning and execution agents, while
remaining accountable for what is claimed. Our thesis is that the binding constraint on
such a team is not the number of agents that can be launched but the researcher's review
attention, and that a framework which makes this constraint explicit also places checks
where the defects that matter occur.

This paper makes the following contributions:
\begin{enumerate}[leftmargin=*]
  \item \textbf{An orchestration framework for research agent teams}
        (Section~\ref{sec:framework}): defined roles, a four-level autonomy scale, and an
        assignment rule that sets each task's level from three properties of its
        effects---reversibility, reach and verifiability. We show that the rule is
        complete and monotone.
  \item \textbf{A file-based handoff protocol} with five statuses and a
        separation-of-acceptance invariant, under which no agent may accept its own
        output, together with a provenance rule for reported numbers that the case
        analysis showed to be necessary.
  \item \textbf{An attention-budget model of review load} that yields the span of control
        of one researcher, shows that the ratio of planning to execution agents follows
        from their review costs rather than being a constant, and relates both to the
        fan-out models of human supervisory control.
  \item \textbf{A three-case evaluation}: an illustrative characterization task from
        analog-to-digital converter (ADC) research, and two real AI-assisted publication
        projects---a 20-chapter textbook and an arXiv paper---whose logs and
        repositories record 13 defects, analyzed by the control that caught each one or
        failed to.
\end{enumerate}

Section~\ref{sec:related} reviews related work on agents, automated science, LLMs in
hardware design, supervisory control and verification. Section~\ref{sec:framework}
presents the framework. Section~\ref{sec:cases} states the research questions and the
case-study method and reports the three cases and the cross-case analysis.
Section~\ref{sec:discussion} discusses implications and limitations, and
Section~\ref{sec:conclusion} concludes with directions for measurement.
